\chapter{Introduction}
\label{ch:introduction}

\section{Motivation: the implementation gap}

In 2024, more than sixty major green hydrogen projects were cancelled globally, removing approximately 4.9 million tonnes per year of would-have-been capacity from public 2030 timelines. BP exited its Australian, Omani, and UK clean hydrogen portfolio. Air Products halted three large US projects. ArcelorMittal walked away from its Bremen and Eisenhüttenstadt plants despite \euro 1.3 billion in committed European Innovation Fund subsidies. In September 2025, seven winners of the EU Hydrogen Bank auction --- representing a combined 1.88 GW of electrolysis capacity --- withdrew from grant negotiations before the security-deposit deadline.\footnote{Sources: industry reporting compiled in \cite{decarbonizeweekly2026,ing2026,bucklebridge2026}; corroborated by S\&P Global Commodity Insights coverage.}

This cancellation wave exposes what \cite{odenweller2025green} term the \emph{green hydrogen implementation gap}: the difference between announced project ambition and realised capacity. Tracking 190 projects over three years, the authors document that only 7\% of 2023 announced capacity reached scheduled Final Investment Decision (FID). Yet the announced pipeline continues to grow --- the global project portfolio nearly tripled to 422 GW over the same period --- creating a widening discrepancy between political ambition and on-the-ground delivery.

The policy response has been substantial. The US Inflation Reduction Act of 2022 introduced a production-tax-credit-style subsidy (Section 45Q) of up to \$85 per tonne of captured CO\textsubscript{2}. The European Union's Innovation Fund disburses over \euro 10 billion in capex grants. The United Kingdom launched the Hydrogen Allocation Round 1 (HAR-1) and Track-1 cluster-tender architecture. China's 14th Five-Year Plan (2021--2025), succeeded by the 15th Five-Year Plan, elevates hydrogen to a strategic priority within national industrial policy. These four policy archetypes --- output credit, capex grant, cluster tender, and state-coordinated mandate --- represent distinct mechanism-design choices, but their relative effectiveness in moving projects from announcement to FID has not been comparatively evaluated.

\section{Research questions}

This thesis addresses three primary research questions:

\begin{enumerate}[label=\textbf{RQ\arabic*.}, leftmargin=*]
\item \textbf{Which carrot-policy mechanisms causally reduce project failure probability, and by how much?} Specifically, do the four policy archetypes identified above produce statistically and economically significant reductions in project cancellation or on-hold status, and how robust are these estimates to alternative identification strategies and parallel-trends violations?

\item \textbf{Through which channels do effective policies operate?} The Dixit-Pindyck (1994) real-options literature suggests two distinct pathways: subsidies can either raise the value-to-investment ratio $V/I$ (making FID more attractive at a fixed uncertainty level) or attack the revenue-volatility component $\sigma$ (lowering the FID threshold itself). Can we empirically distinguish these channels using sector-level heterogeneity in policy effects?

\item \textbf{What is the role of pre-FID offtake commitments?} The 2024--2026 cancellation pattern strongly suggests that subsidies without secured offtake are insufficient. Can we identify a causal effect of pre-FID offtake-commitment on project survival, and does its magnitude approximate or exceed the effects of formal carrot-policies?
\end{enumerate}

A fourth, supplementary question concerns external validity: \textbf{RQ4. What is the quantitative impact of plausible counterfactual policy reforms?} Using estimated treatment effects, we construct five counterfactual scenarios with bootstrap-validated confidence intervals to inform ongoing policy debate.

\section{Contribution}

This thesis advances the literature on hydrogen-policy evaluation in four ways.

\paragraph{Empirical.} It provides the first cross-jurisdictional causal evaluation of carrot-policy mechanisms for low-carbon hydrogen project survival. The four policy comparisons (US 45Q, EU Innovation Fund, UK Track-1, China FYP) are estimated within a unified econometric framework, allowing direct comparison of mechanism effectiveness. Beyond formal policy, it identifies and quantifies the previously under-recognised role of pre-FID offtake commitments.

\paragraph{Methodological.} It combines three complementary identification strategies that have not previously been jointly applied in this field: (i) modern difference-in-differences estimators that address heterogeneous timing (Sun-Abraham 2021; Borusyak-Jaravel-Spiess 2024); (ii) time-varying-parameter difference-in-differences via three state-space specifications (threshold, AR(1), random walk) to detect structural breaks in policy effectiveness; and (iii) Rambachan-Roth (2023) honest sensitivity bounds to assess robustness to parallel-trends violations. The convergence of point estimates across methods, paired with transparent reporting of method-specific robustness, exemplifies what the literature terms \emph{layered transparency} \cite{rambachan2023more}.

\paragraph{Theoretical.} It develops a Dixit-Pindyck (1994) real-options framework that characterises when each carrot-mechanism type dominates, as a function of project-level revenue uncertainty $\sigma$ and value-to-investment ratio $V/I$. The framework yields testable cross-sector predictions --- $\sigma$-attack mechanisms should be more effective in high-volatility sectors (power and heat, transport) than in low-volatility sectors (chemical, refinery) --- which are then validated empirically.

\paragraph{Policy.} It quantifies counterfactual scenarios with concrete numerical implications: an EU-wide hard-link between Innovation Fund eligibility and demonstrable offtake-commitments could deliver an estimated +83 additional FIDs and +858 kt/y of additional green hydrogen capacity over a three-year horizon, and a full sector-optimal carrot mix could add +113 FIDs, +7.83 Mt/y of CO\textsubscript{2} capture, and +1.76 Mt/y of H\textsubscript{2} output.

\section{Strength of claims: what is empirical, interpretive, quasi-causal, and exploratory}
\label{sec:ch1-claims-structure}

The reader will find that the empirical and theoretical contributions of this dissertation occupy several different levels of identification status, and the strength with which any individual finding is asserted varies accordingly. We make this hierarchy explicit at the outset, in line with the disciplined-identification programme of \cite{imbens2018causal} as developed in the Tinbergen-tradition treatment of observational causal inference, and in the framing of \cite{blasques2024mitigating} on \emph{mitigating} estimation risk in observational data rather than eliminating it.

Four categories of claim appear in this dissertation. They are presented in decreasing order of identification strength, and each empirical result reported in subsequent chapters is implicitly or explicitly classified into one of these categories. Appendix~\ref{sec:methods_identification_hierarchy} provides the formal taxonomy; this section provides the reading guide.

\paragraph{(i) Empirically estimated quantities.} Sample averages, regression coefficients with their standard errors, predictive likelihoods, model-confidence-set inclusion statements, out-of-sample forecast loss differentials. These are statistical objects produced by well-specified estimators on the available sample and reported with their associated uncertainty. Examples include the \cite{diebold1995comparing} test statistics for the score-driven specification versus constant-parameter and parameter-driven alternatives in Chapter~\ref{ch:results_tvp}, the matching-estimator point estimates of the offtake-commitment hazard differential in Chapter~\ref{ch:results_offtake}, and the \cite{oster2019unobservable} sensitivity bounds reported throughout. The reader can take these as the most directly identified results in the dissertation.

\paragraph{(ii) Interpretive and organising claims.} Theoretical propositions, mechanism taxonomies, and channel attributions that derive from a model under explicit regularity conditions and that function as organising scaffolding for the empirical results. The five-channel real-options framework of Chapter~\ref{ch:theory} ($\mu$, $\sigma$, $\rho$, $\eta$, $\kappa$) and the credibility-conditional threshold Proposition~\ref{prop:credibility-threshold} are claims of this type. Three formal identification tests reported in Appendix~\ref{app:identification-tests} demonstrate that the channel-attribution claims and the policy-credibility belief are not empirically structurally separable in the observational sample available for this work; the framework therefore functions, in the language we adopt throughout, as a \emph{disciplined interpretive layer} for empirical signatures rather than as a directly identified causal channel decomposition.

\paragraph{(iii) Quasi-causal claims.} Findings produced by identification strategies that approximate but do not achieve full randomised-experiment causal identification: difference-in-differences with parallel-trends robustness, instrumental-variable-style event-study designs around plausibly exogenous policy or macroeconomic discontinuities, propensity-score and inverse-probability-weighted matching estimators with multiple-estimator robustness reporting. The cross-jurisdictional carrot-policy evaluation of Chapter~\ref{ch:results_did}, including the Inflation-Reduction-Act event-window corroboration documented in Paper~2 \cite{saakstra2026paper2}, falls in this category. Quasi-causal claims are reported with their multi-estimator robustness profile and with \cite{rambachan2023more} honest sensitivity bounds where applicable; they are not asserted as definitively structurally identified.

\paragraph{(iv) Exploratory and scenario-based claims.} Counterfactual scenario constructions, hypothetical regime-switching exercises, and policy-mix optimisation calculations. The EU sector-optimal carrot-mix scenario of Chapter~\ref{ch:counterfactual} is reported in this category. These results are constructed as decision-relevant illustrations under explicit assumptions, not as point predictions; their primary contribution is to render the substantive implications of the identified empirical findings legible to a policy-audience reader without overstating their identification status.

The discipline this taxonomy imposes is consequential. Claims at category (i) are asserted directly, with their associated statistical uncertainty. Claims at category (ii) are asserted as theoretically derivable under explicit conditions, with their empirical channel-attribution status reported alongside. Claims at category (iii) are asserted with their multi-estimator robustness profile and identification caveat. Claims at category (iv) are asserted as illustrative under explicit scenario assumptions. The remaining chapters of the dissertation can be read with this hierarchy in mind: when a finding is reported with strong language, the underlying claim category is (i) or (iii) with multi-method convergence; when a finding is reported with cautious language, the category is (ii) or (iv), and the methodological reasons for the caution are documented. Chapter~\ref{ch:discussion} returns to this taxonomy explicitly when discussing the regime-instability findings as substantive contributions rather than as defects of identification.

\section{Structure of the thesis}

Chapter \ref{ch:literature} reviews the existing literature on hydrogen-policy evaluation, real-options approaches to irreversible investment, and modern causal-inference advances in difference-in-differences. Chapter \ref{ch:theory} develops the theoretical framework, formalising the four mechanism types within a real-options model. Chapter \ref{ch:data} describes the project-level S\&P Global dataset, complementary macro variables, and the construction of the analysis panel. Chapter \ref{ch:methodology} presents the econometric methodology, including the rationale for combining static DiD, time-varying DiD, and sensitivity-based identification. Chapters \ref{ch:results_did}--\ref{ch:results_offtake} present the empirical results. Chapter \ref{ch:counterfactual} translates these into counterfactual scenarios. Chapter \ref{ch:discussion} discusses interpretation, limitations, and theoretical implications. Chapter \ref{ch:conclusion} concludes with policy recommendations and avenues for further research.
